\documentclass[11pt]{article}
\usepackage[margin=2.5cm]{geometry}
\usepackage{amsmath,amssymb,amsthm}
\usepackage[hidelinks]{hyperref}
\newcommand{\Q}{\mathbb{Q}}
\newcommand{\Z}{\mathbb{Z}}
\newcommand{\R}{\mathbb{R}}
\newcommand{\p}{\mathfrak{p}}
\newcommand{\m}{\mathfrak{m}}
\newcommand{\code}[1]{\texttt{#1}}
\newtheorem{theorem}{Theorem}
\title{A constructive proof that Bézout does not imply\\ elementary divisor}
\author{Claude (Anthropic)}
\date{6 October 2026}
\begin{document}
\maketitle

\begin{abstract}
Hägg and Mörtberg~\cite{HM}, with GPT-6 Astra, constructed a Bézout domain that is not an elementary
divisor domain; their proof is classical. We outline a constructive proof, based on a Lean formalisation
available at \url{https://github.com/coquand/bezout-constructive}. It has
three ingredients. Chain lifting along explicit polynomial charts replaces compactness.
Smoothness is carried as Jacobian-minor certificates. Points and components of the maximal locus
are explicit. Apart from known constructive results, the only choice principle used is unique
choice.

This note, written by Claude, documents a constructive version of the proof
in~\cite{HM} (\url{https://arxiv.org/abs/2609.35229}) and its Lean formalisation. It has not
been reviewed by Hägg and Mörtberg.
\end{abstract}

\section{The statement}

\begin{theorem}[Hägg--Mörtberg with GPT-6 Astra~\cite{HM}]
There is a Bézout domain $R\supset\Q[x,y]$ over which the matrix
$M=\begin{pmatrix}1+x&y\\y&1-x\end{pmatrix}$ has no diagonal reduction.
\end{theorem}

The ring $R=\bigcup_n A_n$ is an ascending union of smooth, finitely generated $\Q$-domains with
$A_0=\Q[x,y]$. Each step principalises one coprime pair of $A_n$, using a weighted centre
(Abramovich--Temkin--Włodarczyk), the extended Rees algebra and the forcing algebra~\cite{H}
$J_B(c)=B[\sigma]/(\sum c_i\sigma_i-1)$. The original obstruction is topological: compact real sets
with monotone surjections keep the Möbius bundle on $x^2+y^2=1$ nonorientable. The classical proof
uses Noetherianity, maximal ideals, minimal primes and Krull dimension throughout.

\section{Constructive ingredients}

\paragraph{Chain lifting instead of compact sets.}
The topological argument is replaced by a lifting property for finite chains of real points along
each stage. At every stage it is proved with explicit charts and composed along the tower. The
analytic input is a Kantorovich step with explicit radius and a polynomial inverse function
theorem whose radius is uniform in the base point. It uses only positive inverses and positive
roots, with no real division and no Markov principle.

\paragraph{Smoothness as certificates.}
A stage $A=\Q[Y]/(G)$ carries a \emph{certificate of dimension $n$}. This is a comaximal cover by
Jacobian minors $h_k$ such that on each $D(h_k)$ some $m$ of the generators cut out $A$ and their
$m\times m$ minor is invertible. The dimension is $n=|Y|-m$ on every piece. Certificates propagate
along the tower:
\begin{itemize}
\item the forcing algebra $J=B[X_0,\dots,X_r]/(\sum c_iX_i-1)$: on $D(c_i)$ one has
  $J[1/c_i]\cong B[1/c_i][X_j : j\ne i]$, so the $c_i$ themselves give the cover and the
  dimension becomes $n+r$. This replaces the split-conormal argument of \cite[Lemma 3.1]{HM},
  and $B\to J$ is injective iff $\operatorname{Ann}(c)=0$ \cite[Lemma II.2.14]{LQ};
\item the Rees algebra: the weighted charts give the pieces, and the dimension becomes $n+1$.
\end{itemize}
On a piece, $A_{h}$ has the standard presentation
$\Q[Z][Y,t]/(F_{\mathrm{rel}},F_{\mathrm{coord}}-Z,ht-1)$, whose Jacobian is a unit, so $A_h$ is
standard étale over $\Q[Z]$. Formal smoothness at a point follows from Newton lifting. The
formalisation proves the isomorphism by two explicit inverse maps, so no kernel is computed.

\paragraph{Points instead of maximal ideals.}
A \emph{point} of $A$ is a $\Q$-algebra map $A\to L$ onto a finite extension of $\Q$. Its residue
field is therefore discrete. Every local argument of the classical proof is run at a point:
\begin{itemize}
\item the invariant and its admissibility (decided at points);
\item the maximal locus as an intersection over explicit points;
\item primality of $s$ in the Rees algebra;
\item the drop of the invariant at the vertex.
\end{itemize}
Negative goals use $\neg\neg$-stability of membership at a point. The local--global principle for
comaximal localisations \cite{LQ} assembles the pieces. The Nullstellensatz in Jacobson
form produces points: for $\p$ prime and $a\notin\p$ there is a point that kills $\p$ and not $a$.

\paragraph{Explicit inputs.}
\begin{itemize}
\item Factorisation over $\Q$ uses Kronecker's method and Gauss's lemma. It transfers to
  $K(t)$ and to $K[X]/(p)$.
\item Minimal primes follow Perdry's strongly Noetherian approach~\cite{P}.
\item The invariant (order at a point) is the escape level of a Rees bar induction, so no
  unbounded search is needed. The weighted minimiser is a search over an explicit box.
\end{itemize}

\paragraph{Unique choice.}
The tower is a recursively defined sequence of rings, and each stage is selected by a property
that determines it uniquely:
\begin{itemize}
\item the maximal invariant;
\item the canonical component, the least canonical generator list in a decidable order coming
  from an explicit enumeration of the ring;
\item the common denominator, the least admissible index, and the canonical generators of the
  next presentation.
\end{itemize}
Passing from $\exists!$ to data is unique choice. It is valid in Bishop-style mathematics. In
Lean, however, it is an extra principle, because \code{Prop} does not eliminate into \code{Type}.
Nothing stronger is used.


\section{Status of the formalisation}

\begin{itemize}
\item \code{main\_theorem} and \code{main\_theorem\_tower} are proved in Lean~4/Mathlib with the
  standard axioms (9188 build jobs).
\item The constructive layer passes an audit of 1921 declarations. It relies on 8 trusted
  primitives:
  \begin{itemize}
  \item sign splitting and cotransitivity on $\R$;
  \item positive inverses and positive roots;
  \item the Banach fixed point with rate $\tfrac12$;
  \item the polynomial gcd over a discrete field;
  \item the primitive-element shift (a bounded decidable search, justified by $\neg\neg$).
  \end{itemize}
\item In the dependency cone of the main theorem (4812 declarations), no project declaration
  uses classical logic directly. The standard axioms still appear, but only through
  Mathlib lemmas at the level of propositions (Hilbert basis, localisation, finite-set
  bookkeeping), which are known constructive facts. The tower is defined by well-founded
  recursion together with unique choice.
\end{itemize}

\begin{thebibliography}{9}
\bibitem{HM} C.~Hägg and A.~Mörtberg, \emph{A Bézout domain that is not an elementary divisor domain},
  arXiv:2609.35229, 2026, \url{https://arxiv.org/abs/2609.35229}. The proof was found with GPT-6 Astra and formalised in Lean with Codex;
  \url{https://github.com/Zelaron/bezout-counterexample-lean}.
\bibitem{H} M.~Hochster, \emph{Solid closure}, Contemp.\ Math.\ 159 (1994), 103--172.
\bibitem{LQ} H.~Lombardi and C.~Quitté, \emph{Commutative Algebra: Constructive Methods},
  Springer, 2015.
\bibitem{P} H.~Perdry, \emph{Strongly Noetherian rings and constructive ideal theory},
  J.~Symbolic Comput.\ 37 (2004).
\end{thebibliography}

\end{document}
